\documentclass[12pt, a4paper]{article}

% ── Packages ──────────────────────────────────────────────────────────────────
\usepackage[a4paper, margin=2.5cm]{geometry}
\usepackage{graphicx}
\usepackage{float}
\usepackage{amsmath}
\usepackage{amssymb}
\usepackage{xcolor}
\usepackage{mdframed}
\usepackage{parskip}
\usepackage{enumitem}
\usepackage{titlesec}
\usepackage{booktabs}
\usepackage{hyperref}
\usepackage[T1]{fontenc}

% ── Graphics path ─────────────────────────────────────────────────────────────
\graphicspath{{../../Images_for_latex/Lecture_3/}}

% ── Colours ───────────────────────────────────────────────────────────────────
\definecolor{keyblue}{RGB}{30, 80, 160}
\definecolor{keybluebg}{RGB}{220, 235, 255}
\definecolor{noteorange}{RGB}{200, 100, 0}
\definecolor{noteorangebg}{RGB}{255, 240, 215}
\definecolor{warnred}{RGB}{180, 30, 30}
\definecolor{warnredbg}{RGB}{255, 220, 220}
\definecolor{defgreen}{RGB}{20, 110, 50}
\definecolor{defgreenbg}{RGB}{215, 245, 225}
\definecolor{headernavy}{RGB}{25, 35, 90}

% ── Box environments ──────────────────────────────────────────────────────────
\newmdenv[linecolor=keyblue,backgroundcolor=keybluebg,linewidth=1.5pt,
  leftline=true,rightline=false,topline=false,bottomline=false,
  innerleftmargin=10pt,innerrightmargin=8pt,
  innertopmargin=6pt,innerbottommargin=6pt,
  skipabove=8pt,skipbelow=8pt]{keybox}

\newmdenv[linecolor=noteorange,backgroundcolor=noteorangebg,linewidth=1.5pt,
  leftline=true,rightline=false,topline=false,bottomline=false,
  innerleftmargin=10pt,innerrightmargin=8pt,
  innertopmargin=6pt,innerbottommargin=6pt,
  skipabove=8pt,skipbelow=8pt]{notebox}

\newmdenv[linecolor=warnred,backgroundcolor=warnredbg,linewidth=1.5pt,
  leftline=true,rightline=false,topline=false,bottomline=false,
  innerleftmargin=10pt,innerrightmargin=8pt,
  innertopmargin=6pt,innerbottommargin=6pt,
  skipabove=8pt,skipbelow=8pt]{warnbox}

\newmdenv[linecolor=defgreen,backgroundcolor=defgreenbg,linewidth=1.5pt,
  leftline=true,rightline=false,topline=false,bottomline=false,
  innerleftmargin=10pt,innerrightmargin=8pt,
  innertopmargin=6pt,innerbottommargin=6pt,
  skipabove=8pt,skipbelow=8pt]{defbox}

% ── Section formatting ────────────────────────────────────────────────────────
\titleformat{\section}[block]{\large\bfseries\color{headernavy}}{}{0pt}{}
  [\vspace{2pt}\rule{\linewidth}{0.8pt}\vspace{4pt}]
\titleformat{\subsection}[block]{\normalsize\bfseries\color{keyblue}}{}{0pt}{}

% ── Document ──────────────────────────────────────────────────────────────────
\begin{document}

\begin{center}
  {\LARGE\bfseries\color{headernavy} Sensors \& Systems}\\[4pt]
  {\large Kalman Filter for Linear Systems}\\[2pt]
  {\normalsize Torben Knudsen \quad|\quad Electronic Systems, 8th Semester}
\end{center}

\vspace{4pt}\hrule height 1.2pt\vspace{12pt}

% ==============================================================================
\section*{The Three Problem Types: DD, CD and CC}
% ==============================================================================

Every Kalman filter model consists of exactly two equations that
answer two completely different questions:

\begin{itemize}[leftmargin=1.4em]
  \item \textbf{State equation} --- how does the system evolve over
        time?
        This is purely physics: the state changes on its own, driven
        by any input you apply and disturbed by random process noise.
        Nothing is observed yet --- it is just the world doing its
        thing.
  \item \textbf{Measurement equation} --- what does the sensor read
        at a given moment?
        The sensor observes the current state, possibly with a known
        input, and adds its own measurement noise.
        That is all.
\end{itemize}

The choice of \emph{continuous} or \emph{discrete} for the state
equation is a \textbf{modelling decision} based on how you want to
describe the physics.
If the physics are naturally described by differential equations
(chemistry, mechanics) use a continuous state; if you are already
thinking in time steps use a discrete state.
The measurement is almost always discrete in practice because real
sensors sample --- so the measurement equation stays discrete
regardless of what you chose for the state.

Before writing the Kalman filter equations it is important to know
\emph{which version} of the problem you are dealing with.
The type is determined by whether the \textbf{state model} and the
\textbf{measurements} are in discrete or continuous time.

\begin{center}
\renewcommand{\arraystretch}{1.6}
\begin{tabular}{p{1.8cm} p{3cm} p{3cm} p{6.5cm}}
  \toprule
  \textbf{Type} & \textbf{State model} & \textbf{Measurements} &
  \textbf{Typical use} \\
  \midrule
  \textbf{DD} &
    Discrete time & Discrete time &
    Digital systems, embedded processors, sampled data.
    \textbf{This is the focus of this course.} \\
  \textbf{CD} &
    Continuous time (SDE) & Discrete time &
    Physics modelled in continuous time but sensor sampled
    at fixed intervals. \\
  \textbf{CC} &
    Continuous time & Continuous time &
    Theoretical; solution is the Kalman--Bucy filter.
    \emph{Not needed in this course.} \\
  \bottomrule
\end{tabular}
\end{center}
\clearpage
\subsection*{DD --- Discrete-Discrete}

Both the physics model and the sensor are discrete ---
everything happens at fixed sampled steps.

\[
  x_{k+1} = \Phi_k x_k + \Gamma_k u_k + w_k,
  \quad w_k \sim \mathcal{N}(0, Q_k)
\]
\[
  y_k = H_k x_k + D_k u_k + v_k,
  \quad v_k \sim \mathcal{N}(0, R_k)
\]

\begin{center}
\renewcommand{\arraystretch}{1.4}
\begin{tabular}{p{1.5cm} p{12.5cm}}
  \toprule
  \textbf{Symbol} & \textbf{Meaning} \\
  \midrule
  $x_k$ & State vector at step $k$ --- the quantities you want to estimate
          (e.g.\ position, velocity, temperature). \\
  $\Phi_k$ & State transition matrix --- describes how the state naturally
             evolves from step $k$ to step $k{+}1$ with no input. \\
  $\Gamma_k$ & Input matrix --- maps the known control input $u_k$ into
               the state update. \\
  $u_k$ & Known control input at step $k$ (e.g.\ heater power, motor
          voltage). \\
  $w_k$ & Process noise --- random disturbances acting on the state
          (wind, friction, unknown forces). $\sim\mathcal{N}(0,Q_k)$. \\
  $Q_k$ & Process noise covariance --- how large and correlated the
          disturbances are. \\
  $y_k$ & Measurement (sensor reading) at step $k$. \\
  $H_k$ & Measurement matrix --- maps the state to what the sensor
          actually sees (e.g.\ $H=[1\ 0]$ means only position is
          measured, not velocity). \\
  $D_k$ & Direct feedthrough matrix --- the part of the input that
          appears directly in the measurement (often zero). \\
  $v_k$ & Measurement noise --- sensor inaccuracy. $\sim\mathcal{N}(0,R_k)$. \\
  $R_k$ & Measurement noise covariance --- how noisy the sensor is. \\
  \bottomrule
\end{tabular}
\end{center}

The noise sequences $w_k$ and $v_k$ are \textbf{white} (independent
in time) and \textbf{uncorrelated with each other}:
$E(w_k v_l^\top) = 0$.

\begin{warnbox}
  \textbf{Is this a requirement --- what happens if it is violated?}\\
  Yes --- both conditions are \emph{necessary} for the standard
  Kalman filter to be optimal (MMSE).\\[4pt]
  \textbf{If the noise is not white} (i.e.\ it is \emph{coloured}
  --- correlated across time steps, such as a slow drift):
  the filter is no longer optimal because it ignores the memory in
  the noise.
  The fix is to \textbf{augment the state} with a model of the
  coloured noise (e.g.\ model it as a first-order Markov process
  driven by white noise), making the combined noise effectively
  white again.\\[4pt]
  \textbf{If process and measurement noise are correlated}
  ($E(w_k v_k^\top) \neq 0$, e.g.\ both affected by the same
  physical source): the standard equations are wrong.
  A modified version of the KF exists that adds a cross-covariance
  term, but this is uncommon in practice.
\end{warnbox}

\begin{notebox}
  \textbf{Example --- room temperature control.}\\
  The temperature $T_k$ in a room is sampled every second by a
  digital thermometer.
  The heater model is discretised: $T_{k+1} = T_k + \alpha(T_{\text{set}} - T_k)\cdot T_s + w_k$,
  where $w_k$ is random disturbance (window opened, people entering).
  The thermometer reads $y_k = T_k + v_k$.
  Both the state update and the measurement happen at the same
  discrete 1\,s steps --- this is DD.
  The Kalman filter simply multiplies by $\Phi$ to predict and
  corrects with the thermometer reading.
\end{notebox}
\clearpage
\subsection*{CD --- Continuous-Discrete}

The physics happen continuously (described by ODEs/SDEs), but the
sensor only delivers measurements at discrete sample intervals.

\begin{notebox}
  \textbf{ODE vs SDE --- what is the difference?}\\
  An \textbf{ODE} (Ordinary Differential Equation) is a deterministic
  equation describing how a quantity changes over time:
  $\dot{x}(t) = f(x(t),\,t)$.
  Given exact initial conditions it always produces the same
  trajectory --- no randomness.
  Example: Newton's second law $m\ddot{x} = F$ is an ODE.\\[4pt]
  An \textbf{SDE} (Stochastic Differential Equation) is an ODE with
  an added \emph{random noise term}:
  $dx(t) = f(x,t)\,dt + g(x,t)\,dW(t)$,
  where $dW(t)$ is a Wiener process (Brownian motion) increment ---
  continuous-time white noise.
  The same initial conditions produce a \emph{different random
  trajectory} every time.
  SDEs are used when the physics are continuous but genuinely
  uncertain (turbulence, molecular diffusion, financial prices).\\[4pt]
  In the CD Kalman filter the state model is an SDE because real
  continuous systems always have some random disturbance acting on
  them.
\end{notebox}

\[
  dx(t) = \bigl(F(t)x(t) + B(t)u(t)\bigr)dt + d\omega(t),
  \quad \omega(t) \sim \mathcal{W}(Q(t))
\]
\[
  y(t_k) = H_k x(t_k) + D_k u(t_k) + v_k,
  \quad v_k \sim \mathcal{N}(0, R_k)
\]

\begin{center}
\renewcommand{\arraystretch}{1.4}
\begin{tabular}{p{1.8cm} p{12.2cm}}
  \toprule
  \textbf{Symbol} & \textbf{Meaning} \\
  \midrule
  $x(t)$ & State as a \emph{continuous} function of time. \\
  $F(t)$ & Continuous-time system matrix --- equivalent of $\Phi$ but
           for a continuous ODE (not a discrete step). \\
  $B(t)$ & Continuous-time input matrix --- equivalent of $\Gamma$. \\
  $u(t)$ & Continuous control input. \\
  $d\omega(t)$ & Continuous-time process noise increment (Wiener process
                 / Brownian motion) --- the continuous equivalent of $w_k$. \\
  $Q(t)$ & Spectral density of the continuous noise (not the same units
           as discrete $Q_k$). \\
  $y(t_k)$ & Discrete measurement arriving at sample time $t_k$ ---
             same structure as DD. \\
  $H_k,\,D_k,\,v_k,\,R_k$ & Same meaning as in DD, applied at the
                              discrete sample instant $t_k$. \\
  \bottomrule
\end{tabular}
\end{center}

\begin{notebox}
  \textbf{Example --- chemical reactor.}\\
  \begin{itemize}[leftmargin=1.4em, topsep=2pt]
    \item \textbf{Input} $u(t)$: the flow rate of reactant you pump
          into the reactor.
          It is continuous --- you are continuously pumping.
    \item \textbf{State} $x(t)$: the concentration inside the
          reactor.
          It evolves continuously because chemistry does not pause
          between readings:
          $\dot{c}(t) = -k\,c(t) + u(t) + w(t)$.
    \item \textbf{Measurement / output} $y(t_k)$: a lab technician
          scoops out a sample \emph{every 30 minutes} and measures
          the concentration.
          That scoop only happens at discrete moments --- so the
          measurement is discrete even though the concentration
          itself never stopped changing:
          $y_k = c(t_k) + v_k$.
  \end{itemize}
  The state is continuous, the measurement is discrete --- this is
  CD.
\end{notebox}

\begin{warnbox}
  \textbf{Key difference from DD: you must integrate the ODE
  between measurements.}\\
  In DD the prediction step is just $\hat{x}_{k+1|k} =
  \Phi_k \hat{x}_{k|k}$ --- one matrix multiply.\\
  In CD, between two measurement times $t_{k-1}$ and $t_k$ the
  state has been evolving continuously according to the ODE.
  You cannot just multiply by $\Phi$; you must \emph{integrate the
  continuous differential equation forward} from $t_{k-1}$ to $t_k$
  to get the predicted state and covariance.
  For a linear time-invariant system this integral has a closed-form
  solution (matrix exponential); in general it requires a numerical
  ODE solver.
  The measurement update step is then the same as in DD.
\end{warnbox}

\subsection*{CC --- Continuous-Continuous}

Both the state model and the measurements are continuous ---
no sampling anywhere.

In practice CC is almost never implemented directly on hardware
because all modern processors sample.
It exists as the \textbf{theoretical foundation} from which DD and
CD are derived --- understanding CC gives you the intuition, but
for any real digital system you will use DD or CD.

\begin{notebox}
  \textbf{Which type should I use in practice?}\\
  Almost all real implementations run on a digital processor and
  sample sensors at fixed intervals.
  That makes the \textbf{DD Kalman filter} the right choice.
  If the physics are best described in continuous time (e.g.\ a
  mechanical system), you discretise the continuous model first
  (using zero-order hold or Forward Euler) and then apply the DD
  filter.
  The CD filter exists for cases where very accurate discretisation
  of a stiff SDE is needed between measurement steps.
\end{notebox}

% ==============================================================================
\clearpage
\section*{The DD Kalman Filter Algorithm}
% ==============================================================================

This is the complete algorithm --- everything you need to implement
a Kalman filter for a discrete-time linear system.

\subsection*{Notation}

The subscript notation keeps track of \emph{which data was used}
when computing an estimate:

\begin{center}
\renewcommand{\arraystretch}{1.5}
\begin{tabular}{p{3cm} p{5.5cm} p{5.5cm}}
  \toprule
  \textbf{Symbol} & \textbf{Meaning} & \textbf{Plain words} \\
  \midrule
  $\hat{x}_{k|k}$ &
    $E(x_k \mid Y_k)$ &
    Best estimate of $x$ at time $k$, using all data up to and
    including step $k$. \\
  $\hat{x}_{k|k-1}$ &
    $E(x_k \mid Y_{k-1})$ &
    Predicted estimate of $x$ at time $k$, using data only up to
    step $k{-}1$ (before new measurement arrives). \\
  $P_{k|k}$ &
    $E(\tilde{x}_{k|k}\,\tilde{x}_{k|k}^\top)$ &
    Error covariance \emph{after} measurement update. \\
  $P_{k|k-1}$ &
    $E(\tilde{x}_{k|k-1}\,\tilde{x}_{k|k-1}^\top)$ &
    Error covariance \emph{before} measurement (prediction
    uncertainty). \\
  $\tilde{y}_{k|k-1}$ &
    $y_k - \hat{y}_{k|k-1}$ &
    Innovation: how surprising is the new measurement? \\
  \bottomrule
\end{tabular}
\end{center}

\subsection*{Initialisation}
\[
  x_0 \sim \mathcal{N}(\hat{x}_{0|-1},\; P_{0|-1})
\]
Provide a starting guess $\hat{x}_{0|-1}$ for the state and
$P_{0|-1}$ for the uncertainty.
If you have no prior knowledge set $\hat{x}_{0|-1} = 0$ and
$P_{0|-1} = I$.

\subsection*{Step A --- Measurement update (run when $y_k$ arrives)}

A new measurement has just arrived.
Correct the predicted estimate using the new information:

\begin{align}
  \hat{y}_{k|k-1} &= H_k \hat{x}_{k|k-1} + D_k u_k
  \tag{predicted measurement} \\[4pt]
  \tilde{y}_{k|k-1} &= y_k - \hat{y}_{k|k-1}
  \tag{innovation} \\[4pt]
  K_k &= P_{k|k-1} H_k^\top
    \!\bigl(H_k P_{k|k-1} H_k^\top + R_k\bigr)^{-1}
  \tag{Kalman gain} \\[4pt]
  \hat{x}_{k|k} &= \hat{x}_{k|k-1} + K_k \tilde{y}_{k|k-1}
  \tag{updated state estimate} \\[4pt]
  P_{k|k} &= (I - K_k H_k)\,P_{k|k-1}\,(I - K_k H_k)^\top
    + K_k R_k K_k^\top
  \tag{updated covariance}
\end{align}

\subsection*{Step B --- Time update (prediction to next step)}

Propagate the current estimate forward to step $k{+}1$ using the
dynamics model:

\begin{align}
  \hat{x}_{k+1|k} &= \Phi_k \hat{x}_{k|k} + \Gamma_k u_k
  \tag{predicted state} \\[4pt]
  P_{k+1|k} &= \Phi_k P_{k|k} \Phi_k^\top + Q_k
  \tag{predicted covariance}
\end{align}

Then wait for $y_{k+1}$ and repeat Step~A.

\begin{notebox}
  \textbf{What each equation is doing in plain language.}
  \begin{itemize}[leftmargin=1.4em, topsep=2pt]
    \item \textbf{Innovation} $\tilde{y}$: the gap between what the
          sensor actually read and what the filter predicted it would
          read.
          A large gap means the prediction was wrong ---
          trust the measurement.
    \item \textbf{Kalman gain} $K_k$: balances measurement noise
          ($R_k$) against prediction uncertainty ($P_{k|k-1}$).
          Large $R_k$ (noisy sensor) $\Rightarrow$ small $K_k$
          (ignore measurement, trust prediction).
          Large $P_{k|k-1}$ (uncertain prediction) $\Rightarrow$
          large $K_k$ (trust measurement).
    \item \textbf{Updated state}: old prediction nudged towards
          the measurement by fraction $K_k$.
    \item \textbf{Updated covariance}: shrinks because we just
          learned something.
          The $(I - K_k H_k)$ term reduces the uncertainty that
          was already in $P_{k|k-1}$.
    \item \textbf{Predicted covariance}: grows by $Q_k$ because
          the unknown process noise $w_k$ adds fresh uncertainty
          every step.
  \end{itemize}
\end{notebox}

% ==============================================================================
\clearpage
\section*{Block Diagram Interpretation}
% ==============================================================================

The Kalman filter runs two parallel loops --- one for the
\textbf{mean value} (the state estimates) and one for the
\textbf{covariance} (the uncertainties).
They are shown separately below.

\subsection*{Mean value part}

\begin{figure}[H]
  \centering
  \includegraphics[width=\linewidth]{slide-26.jpg}
  \caption{Block diagram of the real system (blue, top) and the
           Kalman filter mean-value part (green, bottom).
           The filter is a copy of the system dynamics driven by
           the same input $u_k$, corrected at each step by the
           Kalman gain $K_k$ applied to the innovation
           $\tilde{y}_{k|k-1}$.}
  \label{fig:slide26}
\end{figure}

The key insight from the diagram: the Kalman filter is essentially
a \textbf{model of the system running in parallel}.
It predicts where the state should be ($\hat{x}_{k|k-1}$ via
$\Phi_k$ and $\Gamma_k$), compares the prediction to the actual
measurement ($\tilde{y}_{k|k-1}$), and adds a correction
proportional to that surprise ($K_k \tilde{y}_{k|k-1}$).
If the model were perfect and there were no noise, the correction
would always be zero.

\begin{notebox}
  \textbf{The filter is a copy of the system.}\\
  Notice the bottom (green) box contains exactly the same blocks as
  the top (blue) system: $\Phi_k$ (state transition), $\Gamma_k$
  (input matrix), $H_k$ (output), and a $q^{-1}$ delay.
  The only additions are $K_k$ (the correction gain) and the
  subtraction node that forms the innovation $\tilde{y}_{k|k-1}
  = y_k - \hat{y}_{k|k-1}$.
  The gain $K_k$ is fed in from the covariance part (dashed arrow).
\end{notebox}

\subsection*{Covariance part}

\begin{figure}[H]
  \centering
  \includegraphics[width=\linewidth]{slide-27.jpg}
  \caption{Block diagram of the Kalman filter covariance part.
           It runs independently of the measurements and feeds the
           gain $K_k$ to the mean-value part.}
  \label{fig:slide27}
\end{figure}

The covariance loop runs \textbf{independently of the measurements}
--- it only uses the model matrices ($\Phi_k$, $H_k$, $Q_k$,
$R_k$), not the actual sensor data $y_k$.
It propagates $P_{k|k-1}$ forward using the time update, then
shrinks it via the measurement update, and outputs the Kalman gain
$K_k$ to the state estimator.

\begin{keybox}
  \textbf{Consequence: the gain can be pre-computed.}\\
  Because $K_k$ and $P_{k|k}$ do not depend on measurements ---
  only on the system model and noise matrices --- for a
  \textbf{time-invariant} system ($\Phi$, $H$, $Q$, $R$ all
  constant) the covariance converges to a steady state $P_\infty$
  and the gain converges to a constant $K_\infty$.
  You can solve for these \emph{offline} before the filter even
  runs, using the \textbf{Discrete Algebraic Riccati Equation
  (DARE)}:
  \[
    P = \Phi P \Phi^\top
      - \Phi P H^\top (H P H^\top + R)^{-1} H P \Phi^\top + Q
  \]
  At steady state the filter becomes a fixed-gain observer --- much
  cheaper to run in real time.
\end{keybox}

% ==============================================================================
\clearpage
\section*{Key Properties of the Kalman Filter}
% ==============================================================================

\subsection*{1 --- Optimal in the MMSE sense}

For a \textbf{linear system with Gaussian noise}, the Kalman filter
is the \textbf{minimum mean-square error (MMSE)} estimator.
No other estimator can produce a lower expected squared error.
This is not a heuristic --- it follows directly from the fact that
the optimal estimate of $x$ given $y$ is the conditional mean
$E(x \mid y)$, and for jointly Gaussian distributions this
conditional mean is exactly what the Kalman filter computes.

\subsection*{2 --- The estimate is unbiased}

\[
  E(x_k - \hat{x}_{k|k}) = 0
\]
On average the filter makes no systematic error.
A biased estimator would consistently over- or under-estimate,
which would be easy to exploit --- so unbiasedness is a basic
requirement for any good estimator.

\subsection*{3 --- Orthogonality principle}

The estimation error is \textbf{uncorrelated with the measurements}:
\[
  E\!\left[(x_k - \hat{x}_{k|k})(y_j - E(y_j))^\top\right] = 0
  \quad \forall j \leq k
\]
In plain words: there is no remaining information in the measurements
that the filter has not already used.
If the error \emph{were} correlated with $y$, you could improve the
estimate by adjusting for it --- so zero correlation is the hallmark
of an optimal estimator.

\subsection*{4 --- Innovation sequence is white}

The innovations $\tilde{y}_{k|k-1}$ at different time steps are
\textbf{mutually uncorrelated}:
\[
  E\!\left[\tilde{y}_{k|k-1}\,\tilde{y}_{j|j-1}^\top\right] = 0
  \quad \text{for } k \neq j
\]

\begin{notebox}
  \textbf{Why this matters for fault detection.}\\
  If the filter model is correct, the innovations should look like
  white noise.
  If the innovations are \emph{correlated} across time steps ---
  e.g.\ consistently positive for several steps in a row --- it
  means either the model is wrong, the noise matrices $Q$ or $R$
  are misspecified, or the sensor has drifted.
  Monitoring the innovation sequence (called \textbf{whiteness
  testing}) is one of the main tools for detecting sensor faults
  and model mismatch in real systems.
\end{notebox}

\subsection*{5 --- Gain and covariance are measurement-independent}

$K_k$, $P_{k|k}$, $P_{k|k-1}$ depend only on the model matrices
$\Phi_k$, $H_k$, $Q_k$, $R_k$ --- not on the actual measurement
values $y_k$.
This means the covariance loop (and the gain schedule) can be
computed offline and stored, making real-time implementation very
efficient.

% ==============================================================================
\clearpage
\section*{Possibilities and Challenges}
% ==============================================================================

\subsection*{What the Kalman filter can be used for}

The Kalman filter is a flexible tool --- the creative part is
choosing what goes into the state vector $x_k$:

\begin{center}
\renewcommand{\arraystretch}{1.5}
\begin{tabular}{p{4cm} p{10cm}}
  \toprule
  \textbf{Application} & \textbf{How it works} \\
  \midrule
  State estimation &
    Classic use: estimate position, velocity, orientation from
    noisy sensor data. \\
  Signal filtering &
    Put the signal model in the state; use KF to smooth noisy
    measurements (equivalent to an optimal Wiener filter). \\
  Parameter estimation &
    Augment the state with unknown parameters
    (e.g.\ bias, stiffness).
    The KF tracks and updates the parameter estimate online. \\
  Multi-step prediction &
    Run the time-update equations forward $j$ steps from
    $\hat{x}_{k|k}$ without a measurement to get
    $\hat{x}_{k+j|k}$ and its covariance. \\
  Fault detection &
    Monitor the innovation $\tilde{y}_{k|k-1}$; if it stops
    looking like white noise a fault is likely present. \\
  Sensor fusion &
    Multiple sensors $\to$ multiple rows in $H_k$; the KF
    optimally combines all of them. \\
  \bottomrule
\end{tabular}
\end{center}

\subsection*{Challenges and limitations}

\begin{itemize}[leftmargin=1.4em]
  \item \textbf{Noise must be white.}
        The KF is optimal only when $w_k$ and $v_k$ are white
        (independent in time).
        \emph{Coloured noise} (e.g.\ low-frequency drift) must be
        handled by augmenting the state with a noise model --- for
        example adding a first-order Markov process driven by white
        noise.

  \item \textbf{Need to know $Q$ and $R$.}
        The noise covariances $Q_k$ and $R_k$ must be specified
        by the designer.
        If they can be measured or estimated from data, great; if
        not, they become tuning parameters.
        Wrong values make the filter suboptimal --- or diverge.

  \item \textbf{Linear systems only.}
        For nonlinear systems the EKF (linearise around current
        estimate) or UKF (sigma points) must be used.
        See Lecture~5 for the IMU application.
\end{itemize}

\begin{keybox}
  \textbf{Summary --- the Kalman filter in one paragraph}\\
  The Kalman filter is the optimal (MMSE) recursive estimator for
  linear Gaussian systems.
  It alternates between a \textbf{prediction step} (propagate
  state and covariance forward using the dynamics model) and a
  \textbf{measurement update step} (correct using the innovation,
  weighted by the Kalman gain).
  The gain is computed from the covariance loop, which runs
  independently of the data and can be pre-solved for
  time-invariant systems via the DARE.
  The filter is unbiased, its error is orthogonal to the
  measurements, and its innovations are white --- any violation
  of these properties signals a model or sensor fault.
\end{keybox}

\end{document}
